\documentclass[10pt,a4paper]{amsart}
\usepackage[margin=19mm]{geometry}
\usepackage{amsmath,amssymb,mathtools}
\usepackage[scr=rsfs,cal=euler]{mathalfa}
\usepackage{xfrac}
\usepackage[unicode,hidelinks,bookmarks=false]{hyperref}
\newcommand{\ie}{i.e.\ }
\newcommand{\cf}{cf.\ }
\newcommand{\resp}{respectively\ }
\newcommand{\Subsection}{\subsection}
\providecommand{\nobreakdash}{\nobreak\hbox{-}}
\setlength{\emergencystretch}{2em}
\multlinegap0pt
\usepackage[shortlabels]{enumitem}
\usepackage{amssymb}
\usepackage{tikz}
\usepackage{float}
\newtheorem{theorem}{Theorem}
\newtheorem{proposition}{Proposition}
\title{General position and mutual-visibility in shadow graphs}
\author{Haritha S., Ullas Chandran S. V}
\date{Source: arXiv:2601.19769v1 (CC BY 4.0)}
\begin{document}
\maketitle

\noindent\emph{Source and licence.} General position and mutual-visibility in shadow graphs. Version arXiv:2601.19769v1 (CC BY 4.0), distributed by arXiv under the Creative Commons Attribution 4.0 International licence, http://creativecommons.org/licenses/by/4.0/, as recorded in the arXiv licence metadata for this exact version and shown on the abstract page https://arxiv.org/abs/2601.19769v1. Full licence notice: \texttt{licenses/arxiv-2601.19769-CC-BY-4.0.txt}.\par\smallskip

\noindent\textit{Editorial scope.} Counted unit: the article's proposition proposition (source0.tex lines 499--506). The article's complete original proof of it is delivered with it (source0.tex lines 507--527, one \texttt{proof} environment of 21 lines). No mathematics is written, repaired or re-derived: the statement and the proof are the article's own lines, in the article's own notation, and no retained line is substituted or re-flowed.  Retained uncounted as the source-local context the proof needs: lines 382--385; lines 470--472; lines 530--537.  Nothing was omitted from any retained span, so the package carries no omission record.  No retained line contains a citation, so the wrapper carries no bibliography and no bibliography entry.  No retained line contains a figure environment, an image-inclusion command or drawing code, so no image is delivered and the package declares no asset dependency.  Editorial formatting only, in addition: an AMS \texttt{amsart} preamble that restates the article's own declarations and macros used by the retained lines, namely the environments \texttt{theorem}, \texttt{proposition} on separate counters, together with the standard packages those lines need.  The retained environments print in this document as Theorem 1, Proposition 1 in order of appearance, whereas the article prints the same items with the numbering of its own full document.  The complete native source of the article as distributed in the arXiv e-print for this exact version is bundled unchanged as \texttt{sources/193485/source0.tex}.
\begin{theorem}\label{thm-bounds}
    For any connected graph $G$,
    $$max\{n(G),2\mu_{i}(G), 2\Delta(G)\} \leq \mu(S(G)) \leq min\{n(G)+\mu(G), 2n(G)-2\}.$$
\end{theorem}

\begin{theorem}\label{mv-tree}
    For any tree $T$ with diameter at least 3, $\mu(S(T)) = n(T)+l(G)$. 
\end{theorem}

\begin{theorem}\label{mv-cycle}
 For a cycle $C_n$, 
 $$ \mu(S(C_n))= \begin{cases}
        6, & \text{if } n=4\\
        n+1, & \text{if } n \in \{3,5,6\}\\
        n, & \text{if } n\geq 7.
    \end{cases}$$
\end{theorem}

\begin{proposition}
    The following holds for any graph $G$,
    \begin{enumerate}[label=\roman*.]
        \item  $\mu(S(G))=2$ if and only if $G \cong P_2$.
      \item  $\mu(S(G))= 4$ if and only if $G$ is either $P_3$ or $C_3$.
    \item There exists no connected graph $G$ for which $\mu(S(G)= 3$ or $\mu(S(G)=5$.
    \end{enumerate}
\end{proposition}

\begin{proof}
  
For part (i), it is straightforward to verify that $\mu(S(G)) = 2$ if $G \cong P_2$. Conversely, suppose that $\mu(S(G)) = 2$. Then the lower bound in Theorem \ref{thm-bounds} implies that $\Delta(G) = 1$, and hence $G \cong P_2$.
 To prove (ii), it is straightforward to verify that 
$\mu(S(P_3)) = \mu(S(C_3)) = 4$. 
Conversely, suppose that $\mu(S(G)) = 4$. 
Then, by the lower bound in Theorem \ref{thm-bounds}, we have 
$\Delta(G) \leq 2$ and $n(G) \leq 4$. 
By part (i), it follows that $\Delta(G) = 2$, and hence $G$ is either a path or a cycle. 
Now, Theorems \ref{mv-cycle} and \ref{mv-tree} together imply that $n(G) = 3$. 
Consequently, $G \cong P_3$ or $G \cong C_3$.

To prove (iii), suppose that there exists a connected graph $G$ such that $\mu(S(G)) = 3$. 
Then, by Theorem \ref{thm-bounds}, we have $\Delta(G) = 1$, which contradicts part (i). 
Similarly, suppose that $\mu(S(G)) = 5$. 
Then Theorem \ref{thm-bounds}, together with part (i), implies that $\Delta(G) = 2$. 
Hence, $G$ is either a path or a cycle. 
However, Theorem \ref{mv-cycle} shows that $G$ must be a path. 
Consequently, by Theorem \ref{mv-tree}, we have $G \cong P_3$, which contradicts part (ii).

\end{proof}

\end{document}
